\documentclass[11pt,a4paper]{article}
\usepackage[margin=1in]{geometry}
\usepackage{amsmath,amssymb,amsthm}
\usepackage{algorithm}
\usepackage{algpseudocode}
\usepackage{booktabs}
\usepackage{hyperref}

\theoremstyle{definition}
\newtheorem{definition}{\textbf{Definition}}[section]
\newtheorem{theorem}{\textbf{Theorem}}[section]
\newtheorem{lemma}[theorem]{\textbf{Lemma}}
\newtheorem{remark}{\textbf{Remark}}[section]

\title{\textbf{One-Shot Pruning for Large Language Models via Activation-Weighted Saliency (Wanda)}}
\author{\textbf{Team Scaibu} \\ \small Email: \texttt{jaydeep.9664920749@gmail.com} \\ \small \textbf{Architecture and Core Engine Team}}
\date{\textbf{October 2026}}

\begin{document}
\maketitle

\section{Definitions and Cognitive Mental Models}

\subsection{Formal Mathematical Definition}
\textbf{Definition (Activation-Weighted Saliency Pruning — Wanda):}
In modern transformer-based Large Language Models (LLMs), post-training retraining or iterative fine-tuning across billions of tokens is prohibitively compute-expensive. Traditional post-training pruning techniques fail catastrophically: unweighted magnitude pruning collapses language generation perplexity at modest 50\% sparsity because LLM activations exhibit extreme emergent \textbf{outlier features} (Dettmers et al., 2022) with magnitudes up to 100 times larger than standard channels.

\textbf{Wanda} (Pruning by \textbf{W}eights \textbf{and} \textbf{A}ctivations) (Sun et al., 2023) performs one-shot post-training pruning without weight retraining by factoring input activation statistics directly into the importance metric. Given a linear weight projection matrix $\mathbf{W} \in \mathbb{R}^{d_{\text{out}} \times d_{\text{in}}}$ and calibration input feature activations $\mathbf{X} \in \mathbb{R}^{N \times d_{\text{in}}}$, the column-wise $L_2$-norm vector of input features is:
{\boldmath
\begin{equation}
\|\mathbf{X}_j\|_2 \triangleq \sqrt{\sum_{k=1}^N X_{k, j}^2}, \quad \forall j \in \{1, \dots, d_{\text{in}}\}
\end{equation}
}
The Wanda saliency score $S_{ij}$ for parameter $W_{ij}$ is defined as the product of absolute weight magnitude and activation column norm:
{\boldmath
\begin{equation}
S_{ij} \triangleq |W_{ij}| \cdot \|\mathbf{X}_j\|_2
\end{equation}
}
Pruning is executed independently on each output row $i \in \{1, \dots, d_{\text{out}}\}$: the lowest $k = \lfloor r \cdot d_{\text{in}} \rfloor$ weights (according to score $S_{ij}$) are set strictly to zero, achieving exact target sparsity ratio $r \in [0, 1)$ without solving expensive second-order Hessian inversions.

\subsection{Expanded Non-Technical Definition (Intuition for Anyone)}
\begin{enumerate}
    \item \textbf{A Tiny Valve on a Giant Firehose}: Imagine a chemical factory with 1,000 pipes and valves. If you only look at valve size (weight magnitude), a tiny valve might look unimportant. But if that tiny valve controls a high-pressure firehose carrying 99\% of the factory's fuel (activation outlier), closing it will blow up the factory! Wanda checks both valve size and fluid flow before pruning.
    \item \textbf{No Expensive Retraining Needed}: Older pruning methods required re-training the AI on massive supercomputers for weeks after cutting connections. Wanda achieves 50\% sparsity in seconds on a single GPU using only a small calibration prompt.
    \item \textbf{Row-Wise Uniformity}: Pruning exactly 50\% of weights in every single row ensures perfectly balanced tensor shapes, preventing hardware execution stalls.
\end{enumerate}

\subsection{Child-Friendly Mental Model (Physical Metaphor)}
Imagine turning off light switches in an office building to save power:
\begin{itemize}
    \item \textbf{Naive Magnitude Pruning}: You unscrew light bulbs that look small, not realizing one tiny bulb is the emergency exit sign.
    \item \textbf{Wanda Pruning}: You check how bright and vital each bulb's circuit is before unscrewing it. You unscrew only bulbs in empty closets where no current is flowing!
\end{itemize}

\section{Core Mathematical Equations and Definitions}

\subsection{Output Activation Reconstruction Error}
{\boldmath
\begin{equation}
\|\Delta \mathbf{y}_i\|_2 = \left\| \sum_{j \in \mathcal{P}_i} W_{ij} \mathbf{X}_j \right\|_2 \le \sum_{j \in \mathcal{P}_i} |W_{ij}| \cdot \|\mathbf{X}_j\|_2 = \sum_{j \in \mathcal{P}_i} S_{ij}
\end{equation}
}
\begin{itemize}
    \item \textbf{Formal Definition}: Upper bound on the output displacement vector $\Delta \mathbf{y}_i$ for row $i$ induced by zeroing pruned subset $\mathcal{P}_i$.
    \item \textbf{What It Means}: Minimizing the sum of Wanda scores $\sum S_{ij}$ directly minimizes the worst-case distortion of the transformer block's output activations.
    \item \textbf{Why It Is Important}: Explains why Wanda retains near-dense perplexity without requiring fine-tuning.
\end{itemize}

\subsection{Sparsity and Parameter Zeroing Conservation}
{\boldmath
\begin{equation}
\mathcal{S}_{\text{achieved}} = \frac{1}{d_{\text{out}} \cdot d_{\text{in}}} \sum_{i=1}^{d_{\text{out}}} \sum_{j=1}^{d_{\text{in}}} \mathbf{1}\left( W_{ij}^{\text{pruned}} = 0 \right) \equiv \frac{\lfloor r \cdot d_{\text{in}} \rfloor}{d_{\text{in}}}
\end{equation}
}
\begin{itemize}
    \item \textbf{Formal Definition}: Exact fractional sparsity metric conserved identically across all matrix rows.
    \item \textbf{What It Means}: Guarantees uniform density across all attention projection layers ($W_q, W_k, W_v, W_o$) and MLP feed-forward layers ($W_{\text{gate}}, W_{\text{up}}, W_{\text{down}}$).
    \item \textbf{Why It Is Important}: Enables consistent memory bandwidth reduction during KV cache and matrix multiplication kernels.
\end{itemize}

\section{Rigorous Mathematical Analysis Checklist}

\begin{itemize}
    \item \textbf{Notation}: $\mathbf{W} \in \mathbb{R}^{d_{\text{out}} \times d_{\text{in}}}$, $\mathbf{X} \in \mathbb{R}^{N \times d_{\text{in}}}$, $S_{ij} \in \mathbb{R}_+$ (saliency), $r \in [0, 1)$ (pruning ratio), $\mathcal{P}_i$ (pruned column set).
    \item \textbf{Epistemic Status}: Validated across all leading open LLM architectures (Llama 2/3, Mistral, Gemma).
    \item \textbf{Dependency Graph}: Calibration Input $\to$ Column Norm $\|\mathbf{X}_j\|_2 \to$ Saliency Product $S_{ij} \to$ Row-Wise Top-$K$ Selection $\to$ Zeroed Matrix.
    \item \textbf{Dimensions}: Weight matrix $\mathbb{R}^{d_{\text{out}} \times d_{\text{in}}}$; activation norm $\mathbb{R}^{d_{\text{in}}}$; score matrix $\mathbb{R}^{d_{\text{out}} \times d_{\text{in}}}$.
    \item \textbf{Regimes}: One-shot compression of LLMs from 7B to 70B+ parameter scales.
    \item \textbf{Invariants}: Exact row-wise sparsity invariance: each row has exactly $\lfloor r \cdot d_{\text{in}} \rfloor$ zeroed entries.
    \item \textbf{Isomorphisms}: First-order diagonal approximation to Optimal Brain Surgeon (OBS) and SparseGPT inversion.
\end{itemize}

\section{Theoretical Theorems and Formal Proofs}

\begin{theorem}[Minimization of Output Layer Perturbation Bound]
For a linear layer $\mathbf{y}_i = \sum_{j=1}^{d_{\text{in}}} W_{ij} \mathbf{x}_j$, let $\mathcal{P}_i \subset \{1, \dots, d_{\text{in}}\}$ be the set of pruned column indices with cardinality $|\mathcal{P}_i| = k$. The Wanda pruning strategy that selects $\mathcal{P}_i^* = \operatorname{argmin}_{|\mathcal{P}_i|=k} \sum_{j \in \mathcal{P}_i} |W_{ij}| \cdot \|\mathbf{x}_j\|_2$ minimizes the uniform Cauchy-Schwarz upper bound on the output error magnitude $\|\Delta \mathbf{y}_i\|_2$:
{\boldmath
\begin{equation}
\left\| \Delta \mathbf{y}_i \right\|_2 \le \sum_{j \in \mathcal{P}_i^*} S_{ij} \le \sum_{j \in \mathcal{P}_i} S_{ij}, \quad \forall \, \mathcal{P}_i \text{ s.t. } |\mathcal{P}_i| = k
\end{equation}
}
\end{theorem}

\begin{proof}
The output error vector $\Delta \mathbf{y}_i$ induced by setting weights in $\mathcal{P}_i$ to zero is:
\begin{equation}
\Delta \mathbf{y}_i = \mathbf{y}_i - \mathbf{y}_i^{\text{pruned}} = \sum_{j \in \mathcal{P}_i} W_{ij} \mathbf{x}_j
\end{equation}
Taking the Euclidean norm and applying the triangle inequality:
\begin{equation}
\|\Delta \mathbf{y}_i\|_2 = \left\| \sum_{j \in \mathcal{P}_i} W_{ij} \mathbf{x}_j \right\|_2 \le \sum_{j \in \mathcal{P}_i} \|W_{ij} \mathbf{x}_j\|_2 = \sum_{j \in \mathcal{P}_i} |W_{ij}| \cdot \|\mathbf{x}_j\|_2 = \sum_{j \in \mathcal{P}_i} S_{ij}
\end{equation}
To minimize this upper bound subject to the cardinality constraint $|\mathcal{P}_i| = k$, we solve:
\begin{equation}
\min_{\mathcal{P}_i \subset \{1, \dots, d_{\text{in}}\}, \, |\mathcal{P}_i|=k} \sum_{j \in \mathcal{P}_i} S_{ij}
\end{equation}
Because the sum is separable across coordinates $j$, the global minimum is achieved greedily by selecting the $k$ smallest values of $S_{ij}$.
Hence, Wanda's selection rule produces the optimal set $\mathcal{P}_i^*$ that minimizes the upper bound on layer reconstruction distortion.
\end{proof}

\begin{theorem}[Failure of Pure Magnitude Pruning Under Activation Outliers]
Assume input channel $k$ contains an emergent activation outlier feature such that $\|\mathbf{X}_k\|_2 = M \gg 1$, while for all other channels $j \neq k$, $\|\mathbf{X}_j\|_2 \le 1$. Suppose weight $|W_{ik}| = \epsilon > 0$ is small such that $|W_{ik}| < |W_{ij}|$ for all $j \neq k$, but $\epsilon M > \max_j |W_{ij}|$. Then:
\begin{enumerate}
    \item Pure magnitude pruning prunes coordinate $k$, incurring an error bound $\Delta_{\text{mag}} = \epsilon M \gg 1$.
    \item Wanda preserves coordinate $k$, selecting coordinate $j^*$ with error $\Delta_{\text{wanda}} = |W_{ij^*}| \le \max_j |W_{ij}| \ll \epsilon M$.
\end{enumerate}
\end{theorem}

\begin{proof}
Under pure magnitude pruning, the score function is $S_{ij}^{\text{mag}} = |W_{ij}|$.
Because $|W_{ik}| = \epsilon < |W_{ij}|$ for all $j \neq k$, coordinate $k$ has the strictly lowest magnitude and is pruned first. The resulting output error contribution is:
\begin{equation}
\Delta_{\text{mag}} = |W_{ik}| \cdot \|\mathbf{X}_k\|_2 = \epsilon M
\end{equation}
Under Wanda, the score function is $S_{ij}^{\text{wanda}} = |W_{ij}| \cdot \|\mathbf{X}_j\|_2$.
For the outlier coordinate $k$:
\begin{equation}
S_{ik}^{\text{wanda}} = |W_{ik}| \cdot \|\mathbf{X}_k\|_2 = \epsilon M
\end{equation}
For any other coordinate $j \neq k$:
\begin{equation}
S_{ij}^{\text{wanda}} = |W_{ij}| \cdot \|\mathbf{X}_j\|_2 \le |W_{ij}| \le \max_m |W_{im}|
\end{equation}
By premise, $\epsilon M > \max_m |W_{im}|$. Therefore:
\begin{equation}
S_{ik}^{\text{wanda}} > S_{ij}^{\text{wanda}}, \quad \forall j \neq k
\end{equation}
Consequently, coordinate $k$ will never be in the lowest $k$ scores and will be preserved by Wanda. Wanda instead prunes a small non-outlier coordinate $j^*$, incurring an error bound of at most $\max_m |W_{im}| \ll \epsilon M$. This proves that Wanda prevents catastrophic perplexity blowup caused by destroying outlier activation channels.
\end{proof}

\section{Architectural Questions and Real-World Trade-Offs}

\subsection{Architectural Question 1: Calibration Dataset Sensitivity}
\textbf{Question:} How many calibration sequences are needed to compute stable activation norms $\|\mathbf{X}_j\|_2$?\\
\textbf{Answer:} Remarkably few. Empirical benchmarks on Llama 2/3 demonstrate that 128 sequences of 2,048 tokens drawn from C4 or WikiText are sufficient. Because activation outliers are structural properties of the model's learned representation rather than content-specific anomalies, the column norms converge rapidly, and additional calibration data produces negligible change in the resulting sparsity mask.

\subsection{Architectural Question 2: Comparison with SparseGPT (Hessian Inversion)}
\textbf{Question:} Why choose Wanda over second-order Hessian methods like SparseGPT?\\
\textbf{Answer:} SparseGPT computes the inverse Hessian $\mathbf{H}^{-1} = (\mathbf{X}^\top \mathbf{X} + \lambda \mathbf{I})^{-1}$ to perform weight compensation after each pruned column. While SparseGPT achieves slightly lower perplexity at high sparsity ($>60\%$), computing and factorizing $\mathbf{H}^{-1}$ takes substantial time and can suffer from numerical Cholesky instability in FP16. Wanda requires zero matrix inversions, zero weight updates, and runs $>10\times$ faster, making it ideal for rapid post-training quantization and pruning workflows.

\section{Algorithmic Implementation Specification}

\begin{algorithm}[H]
\caption{One-Shot Structured Pruning for LLMs (Wanda)}
\begin{algorithmic}[1]
\Require Weight matrix $\mathbf{W} \in \mathbb{R}^{d_{\text{out}} \times d_{\text{in}}}$, activation norms $\mathbf{a} \in \mathbb{R}^{d_{\text{in}}}$, ratio $r \in [0, 1)$, mode $\in \{\text{"wanda"}, \text{"magnitude"}\}$
\Ensure Pruned matrix $\mathbf{W}' \in \mathbb{R}^{d_{\text{out}} \times d_{\text{in}}}$, score matrix $\mathbf{S}$, achieved sparsity $s_{\text{ach}}$, zeroed count $N_{\text{zero}}$
\State $d_{\text{out}} \gets \operatorname{len}(\mathbf{W}), \quad d_{\text{in}} \gets \operatorname{len}(\mathbf{a})$
\State $k_{\text{prune}} \gets \lfloor r \cdot d_{\text{in}} \rfloor$
\State Initialize $\mathbf{W}' \gets [\,], \quad \mathbf{S} \gets [\,], \quad N_{\text{zero}} \gets 0$
\For{$i = 0$ \textbf{to} $d_{\text{out}} - 1$}
    \State Initialize row scores $\mathbf{s}_i \gets [\,]$
    \For{$j = 0$ \textbf{to} $d_{\text{in}} - 1$}
        \If{mode $=$ "magnitude"}
            \State Append $|W_{ij}|$ to $\mathbf{s}_i$
        \Else
            \State Append $|W_{ij}| \cdot a_j$ to $\mathbf{s}_i$
        \EndIf
    \EndFor
    \State Append $\mathbf{s}_i$ to $\mathbf{S}$
    \If{$k_{\text{prune}} > 0$}
        \State $\boldsymbol{\pi}_i \gets \operatorname{argsort}(\mathbf{s}_i)$
        \State $\mathcal{P}_i \gets \operatorname{Set}(\boldsymbol{\pi}_i[0 : k_{\text{prune}}])$
        \State Append $[ 0.0 \text{ if } j \in \mathcal{P}_i \text{ else } W_{ij} \text{ for } j \in [0, d_{\text{in}}-1] ]$ to $\mathbf{W}'$
        \State $N_{\text{zero}} \gets N_{\text{zero}} + k_{\text{prune}}$
    \Else
        \State Append $\mathbf{W}[i]$ to $\mathbf{W}'$
    \EndIf
\EndFor
\State $s_{\text{ach}} \gets N_{\text{zero}} / (d_{\text{out}} \cdot d_{\text{in}})$
\State \Return $(\mathbf{W}', \mathbf{S}, s_{\text{ach}}, N_{\text{zero}}, d_{\text{out}} \cdot d_{\text{in}})$
\end{algorithmic}
\end{algorithm}

\section{Big-$\Theta$ Complexity Analysis}

\begin{itemize}
    \item \textbf{Score Evaluation Pass}: $\Theta(d_{\text{out}} \cdot d_{\text{in}})$ scalar multiplications.
    \item \textbf{Row-Wise Sorting}: $\Theta(d_{\text{out}} \cdot d_{\text{in}} \ln d_{\text{in}})$ column index sort.
    \item \textbf{Matrix Pruning Application}: $\Theta(d_{\text{out}} \cdot d_{\text{in}})$ linear assignment pass.
    \item \textbf{Total Time Complexity}: $\Theta(d_{\text{out}} \cdot d_{\text{in}} \ln d_{\text{in}})$ bounded by row sorting.
    \item \textbf{Space Complexity}: $\Theta(d_{\text{out}} \cdot d_{\text{in}})$ storage for saliency scores and output matrix.
\end{itemize}

\section{Single Canonical Repository Reference}

The complete deterministic scalar implementation, verification suite, and unit test benchmarks are published at:
\begin{center}
\url{https://github.com/Chief-Strategist-J/llm-obs-infra}
\end{center}

\end{document}
